\chapter{Elective Caesarean Section}

\section*{Definition and Overview}
An elective caesarean section is a planned surgical procedure in which a baby is delivered through an incision in the mother's abdomen and uterus, rather than through the vagina. It is arranged before labour begins, usually at around 39 weeks of pregnancy, and may be recommended for medical reasons or chosen by the woman. An elective caesarean is a major operation, and while it is generally safe, it carries more risks than an uncomplicated vaginal birth. This chapter explains when an elective caesarean is recommended, what the operation involves, and the recovery and considerations afterwards.

\section*{Epidemiology}
The caesarean rate in Australia is around one in three births, and a large proportion of these are elective. Common reasons include a previous caesarean, breech presentation, placenta praevia, and maternal medical conditions, and an increasing number of women choose a caesarean in the absence of a medical indication. The operation is safe in high-resource settings, but it is associated with higher rates of some complications than vaginal birth, including infection, bleeding, and problems in future pregnancies.

\section*{Aetiology and Pathophysiology}
Elective caesarean section is performed for a range of indications.
\begin{itemize}
    \item \textbf{Previous caesarean:} many women plan a repeat caesarean, while others attempt vaginal birth after caesarean
    \item \textbf{Breech presentation:} a baby positioned feet or bottom first at term is often delivered by caesarean
    \item \textbf{Placenta praevia:} a placenta covering the cervix prevents vaginal birth
    \item \textbf{Maternal medical conditions:} such as severe pre-eclampsia, heart disease, or diabetes with complications
    \item \textbf{Fetal conditions:} growth restriction, fetal anomalies, or large babies
    \item \textbf{Multiple pregnancy:} twins or higher multiples, depending on presentation
    \item \textbf{Maternal choice:} a caesarean requested without a medical indication after full counselling
    \item \textbf{The procedure:} a horizontal incision is made low on the abdomen and uterus, the baby and placenta are delivered, and the layers are closed
\end{itemize}

\section*{Clinical Features}
An elective caesarean is planned and involves preparation.
\begin{itemize}
    \item Discussion and informed consent, including the reasons and the risks and benefits
    \item Blood tests and pre-operative assessment before the operation
    \item Fasting from food and drink for a period before surgery
    \item Admission on the day, with the operation scheduled at around 39 weeks
    \item A spinal or epidural anaesthetic, so the mother is awake for the birth
    \item The operation itself takes about 45 minutes to an hour
    \item The baby is checked immediately after birth, with skin-to-skin contact encouraged
\end{itemize}

\section*{Differential Diagnosis}
The decision between caesarean and vaginal birth is individualised.
\begin{itemize}
    \item Vaginal birth, which is usually recommended when there is no medical reason for caesarean
    \item Vaginal birth after a previous caesarean, which is a safe option for many women
    \item Emergency caesarean, which differs from an elective procedure
    \item Induction of labour as an alternative in some circumstances
\end{itemize}

\section*{When to Seek Medical Care}
Women planning a caesarean should attend all antenatal appointments and pre-operative reviews. They should seek immediate care if labour starts before the planned date, if waters break, or if there is bleeding, pain, or reduced fetal movement.

\textbf{Call 000 (or local emergency services) for:}
\begin{itemize}
    \item Signs of labour before the planned caesarean date
    \item Heavy vaginal bleeding
    \item Severe abdominal pain
    \item Reduced fetal movements
    \item After surgery: severe pain, heavy bleeding, fever, or signs of wound infection
\end{itemize}

\section*{Diagnosis and Investigations}
Before an elective caesarean, assessment ensures the operation is as safe as possible.
\begin{itemize}
    \item \textbf{Pre-operative assessment:} blood pressure, blood tests, and review of general health
    \item \textbf{Ultrasound:} confirming the baby's position, size, and the location of the placenta
    \item \textbf{Anaesthetic review:} discussion of the anaesthetic options
    \item \textbf{Planning:} timing at around 39 weeks to balance maturity against the risk of labour starting
    \item \textbf{Blood group and crossmatch:} in case a transfusion is needed
\end{itemize}

\section*{Management}
An elective caesarean involves the whole maternity team.
\begin{itemize}
    \item \textbf{Anaesthesia:} spinal or epidural, allowing the mother to be awake, with general anaesthesia used when needed
    \item \textbf{The operation:} a low transverse incision in the abdomen and uterus
    \item \textbf{Delivery:} the baby is delivered, and the cord is clamped after a brief delay where possible
    \item \textbf{Closure:} the uterus and abdominal layers are repaired
    \item \textbf{Post-operative care:} pain relief, early mobilisation, wound care, and support with breastfeeding
    \item \textbf{Hospital stay:} usually two to four days
    \item \textbf{Recovery at home:} wound care, activity limits, and follow-up with the GP and obstetrician
\end{itemize}

\section*{Complications}
\begin{itemize}
    \item Infection of the wound, uterus, or bladder
    \item Heavy bleeding requiring transfusion
    \item Blood clots in the legs or lungs
    \item Damage to surrounding organs in rare cases
    \item Complications of anaesthesia
    \item Longer recovery and more post-operative pain than after vaginal birth
    \item Effects on future pregnancies, including a higher risk of complications and the usual recommendation for repeat caesarean
\end{itemize}

\section*{Prognosis and Outcomes}
The outlook after an elective caesarean is very good. The operation is safe, and most women recover well within six weeks and return to full activity. Breastfeeding, bonding, and recovery are supported with good care. The main long-term considerations relate to future pregnancies, which are planned with the obstetric team.

\section*{Prevention and Self-Care}
\begin{itemize}
    \item Attend all antenatal and pre-operative appointments
    \item Ask questions and understand the reasons for and risks of the operation
    \item Follow fasting and preparation instructions
    \item After surgery: take prescribed pain relief, mobilise early, and care for the wound
    \item Avoid heavy lifting and driving until cleared
    \item Seek help with breastfeeding and newborn care
    \item Attend the postnatal and six-week reviews
\end{itemize}

\section*{Sources and Further Reading}
The following reputable sources provide further information for healthcare students and patients seeking reliable, current advice about elective caesarean section.
\begin{itemize}
    \item Healthdirect Australia. \textit{Caesarean section.} https://www.healthdirect.gov.au/caesarean-section
    \item Royal Australian and New Zealand College of Obstetricians and Gynaecologists. \textit{Caesarean section information for women.} https://ranzcog.edu.au/
    \item NHS. \textit{Caesarean section.} https://www.nhs.uk/conditions/caesarean-section/
    \item Mayo Clinic. \textit{C-section: what to expect.} https://www.mayoclinic.org/tests-procedures/c-section/
    \item The Royal Women's Hospital. \textit{Caesarean section information.} https://www.thewomens.org.au/
    \item MedlinePlus. \textit{Cesarean section.} https://medlineplus.gov/cesareansection.html
\end{itemize}
